\thispagestyle{empty}
\vspace*{\fill}
\begin{center}
{\sffamily\bfseries\Huge\color{tmNavy} Matem\'aticas para\\Ingenier\'ia de Sistemas}\\[10pt]
{\color{tmOro}\rule{6cm}{1.4pt}}\\[10pt]
{\sffamily\large Gisselle Severich \;\textbullet\; Luzmila Pati \;\textbullet\; Jery Loza \;\textbullet\; Isacc Mealla}\\[6pt]
{\sffamily\small\color{tmGris} Nombre de la editorial}
\end{center}
\vspace*{\fill}
\cleardoublepage

\thispagestyle{empty}
\vspace*{\fill}
{\sffamily\small\color{tmGris}
\textbf{Matem\'aticas para Ingenier\'ia de Sistemas}\\
Gisselle Severich \;\textbullet\; Luzmila Pati \;\textbullet\; Jery Loza \;\textbullet\; Isacc Mealla\\[4pt]
Primera edici\'on, 2026.\\
Nombre de la editorial.\\[4pt]
Todos los derechos reservados. Queda prohibida la reproducci\'on total o parcial de esta obra sin autorizaci\'on escrita de los autores.\\[4pt]
Los retos interactivos de este libro funcionan en la plataforma TortuMath:\\
\texttt{https://isaccmealla.github.io/TortuMath/}}
\cleardoublepage

\chapter*{C\'omo usar este libro}
\markboth{C\'omo usar este libro}{}
\addcontentsline{toc}{chapter}{C\'omo usar este libro}
\tmgancho{\'Este no es un libro para leer de corrido: es un libro para \textbf{hacer}. Cada secci\'on es corta, visual y termina en un reto.}

Todas las secciones siguen la misma receta. As\'i nunca te pierdes:

\begin{tmdef}{Definici\'on (caja azul)}
Aqu\'i va la idea central y su f\'ormula, en m\'aximo tres l\'ineas.
\end{tmdef}

\begin{tmejemplo}{Ejemplo resuelto (caja verde)}
\begin{tmpasos}
\item Un paso corto y claro.
\item Otro paso, siempre con su \textbf{comprobaci\'on} al final.
\end{tmpasos}
\end{tmejemplo}

\begin{tmsistemas}
Esta caja \'ambar te muestra d\'onde aparece la idea en el trabajo real de un ingeniero de sistemas.
\end{tmsistemas}

\begin{tmtip}Los trucos te ahorran tiempo en los ex\'amenes.\end{tmtip}
\begin{tmcuidado}Aqu\'i viven los errores m\'as comunes. L\'eelos antes de equivocarte.\end{tmcuidado}

\subsection{Niveles de los ejercicios}
\begin{center}
\begin{tabular}{@{}cl@{\hspace{2cm}}cl@{\hspace{2cm}}cl@{}}
\tmbasico & B\'asico & \tmmedio & Medio & \tmdificil & Reto
\end{tabular}
\end{center}

\subsection{Los retos TortuMath}
Al final de cada secci\'on encontrar\'as un recuadro verde con un c\'odigo QR. Escan\'ealo con tu celular o abre el enlace: te llevar\'a al laboratorio interactivo del cap\'itulo, donde resuelves tres misiones y acreditas \textbf{TortuPuntos}.

\tmreto{1}{Prueba el formato: entra al laboratorio del cap\'itulo 1 y resuelve la misi\'on de la zona segura.}

\subsection{Plan de estudio sugerido (un semestre)}
\begin{center}
\begin{tabular}{|c|p{4.6cm}|p{6.4cm}|}
\hline
\rowcolor{tmAzulClaro}\textbf{Semanas} & \textbf{Contenido} & \textbf{Qu\'e entregas}\\ \hline
1 a 3 & Cap\'itulo 1: fundamentos & Ejercicios 1.1 a 1.5 y reto 1\\ \hline
4 a 6 & Cap\'itulo 2: discreta & Ejercicios 2.1 a 2.7 y reto 2\\ \hline
7 a 9 & Cap\'itulo 3: \'algebra & Ejercicios 3.1 a 3.4 y reto 3\\ \hline
10 a 12 & Cap\'itulo 4: c\'alculo & Ejercicios 4.1 a 4.5 y reto 4\\ \hline
13 a 15 & Cap\'itulo 5: probabilidad & Ejercicios 5.1 a 5.5 y reto 5\\ \hline
16 & Proyecto final & Entrega y presentaci\'on (Ap\'endice C)\\ \hline
\end{tabular}
\end{center}

\subsection{C\'omo estudiar cada secci\'on}
\begin{tmpasos}
\item Lee el gancho y la definici\'on (2 minutos).
\item Reproduce a mano un ejemplo resuelto sin mirar los pasos.
\item Resuelve los ejercicios \tmbasico\ y \tmmedio; deja los \tmdificil\ para el final.
\item Haz el reto TortuMath y revisa las respuestas del Ap\'endice A.
\item Cierra con el resumen y, al terminar el cap\'itulo, la autoevaluaci\'on.
\end{tmpasos}

\begin{tmtip}Estudiar en sesiones de 25 minutos con descansos de 5 rinde m\'as que una sola sesi\'on larga. Un cap\'itulo de este libro se estudia bien en tres de esas sesiones.\end{tmtip}

\subsection{Mapa del libro}
\begin{center}
\begin{tikzpicture}[every node/.style={font=\sffamily\small}]
  \node[circle,fill=tmAzul,text=white,minimum size=1.1cm,font=\sffamily\bfseries\Large] (a) at (0,0) {1};
  \node[circle,fill=tmMorado,text=white,minimum size=1.1cm,font=\sffamily\bfseries\Large] (b) at (2.7,0) {2};
  \node[circle,fill=tmVerde,text=white,minimum size=1.1cm,font=\sffamily\bfseries\Large] (c) at (5.4,0) {3};
  \node[circle,fill=tmAmbar,text=white,minimum size=1.1cm,font=\sffamily\bfseries\Large] (d) at (8.1,0) {4};
  \node[circle,fill=tmCoral,text=white,minimum size=1.1cm,font=\sffamily\bfseries\Large] (e) at (10.8,0) {5};
  \node[text=tmNavy,font=\sffamily\small\bfseries] at (0,-1.1) {Fundamentos};
  \node[text=tmNavy,font=\sffamily\small\bfseries] at (2.7,-1.1) {Discreta};
  \node[text=tmNavy,font=\sffamily\small\bfseries] at (5.4,-1.1) {\'Algebra};
  \node[text=tmNavy,font=\sffamily\small\bfseries] at (8.1,-1.1) {C\'alculo};
  \node[text=tmNavy,font=\sffamily\small\bfseries] at (10.8,-1.1) {Probabilidad};
  \draw[tmOro,line width=1.4pt,-{Stealth}] (a) -- (b);
  \draw[tmOro,line width=1.4pt,-{Stealth}] (b) -- (c);
  \draw[tmOro,line width=1.4pt,-{Stealth}] (c) -- (d);
  \draw[tmOro,line width=1.4pt,-{Stealth}] (d) -- (e);
\end{tikzpicture}
\end{center}
